\section{Related work}\label{sec:related}

\subsection{Recurrent and continuous-time neural networks}

Recurrent networks such as the LSTM \cite{hochreiter1997lstm} and the GRU
\cite{cho2014gru} are the standard neural models for sequences. They
update their state once per observation and have no notion of the time
between observations; the usual fix is to give the time gap as an extra
input, as we do for our baselines.

Neural ordinary differential equations \cite{chen2018neuralode} define a
hidden state that evolves continuously in time, which makes them a natural
fit for irregular data. Their solution, however, is fixed by the initial
condition, so later observations must be included by other means.
Neural controlled differential equations \cite{kidger2020ncde} solve this
by letting the interpolated data path drive the hidden state, as in
Section~\ref{sec:neural}. Kidger et al.\ evaluated them on the
CharacterTrajectories dataset of the UEA archive with 30, 50 and 70\% of
the observations dropped, chosen uniformly at random and independently
for each series, the same for every model. With five runs, the Neural CDE
reached 98.7\%, 98.8\% and 98.6\% test accuracy, against 93.6\%, 91.3\%
and 90.4\% for a GRU with the time gap as input, and above other ODE- and
RNN-based models such as GRU-D and ODE-RNN \cite{kidger2020ncde}. Our
missing-data protocol follows theirs. For long series, Neural rough
differential equations \cite{morrill2021nrde} replace the path by
log-signatures over short intervals, which reduces the number of solver
steps.

Liquid time-constant networks \cite{hasani2021ltc} build each neuron as a
linear first-order dynamical system whose time constant is modulated by
nonlinear gates; the authors show that these dynamics are stable and
bounded and report improved results on time-series prediction tasks
compared with classical and modern RNNs. Because the state is computed by
a numerical ODE solver, LTCs are slow. Closed-form continuous-time networks
\cite{hasani2022cfc} approximate the solution of the LTC dynamics in
closed form, with the time appearing explicitly in the formula; the
authors report models that are between one and five orders of magnitude
faster in training and inference than their ODE-based counterparts. Our
case study measures this difference in cost directly
(Table~\ref{tab:cost}).

\subsection{Path signatures}

The signature goes back to Chen \cite{chen1958integration}, who showed
that the iterated integrals of a path combine by what is now called Chen's
identity. Hambly and Lyons \cite{hambly2010uniqueness} proved that the
signature determines a path of bounded variation up to tree-like
equivalence, which makes it a faithful summary of a series. Introductions
for machine learning are given by Chevyrev and Kormilitzin
\cite{chevyrev2026primer} and by Lyons and McLeod \cite{lyons2022sigml}.

Morrill et al.\ \cite{morrill2020generalised} unify the many variants of
the signature method (augmentations such as the time channel and the
basepoint, windows, and the choice of signature or log-signature) into
one generalised signature method. From a study on 26 datasets they derive
a default set of choices and show that it is competitive with benchmark
methods for multivariate time series classification; our signature
features use two of these augmentations. Signatures have also been used
inside neural networks, as a layer that can be trained
\cite{kidger2019deepsig}, to view recurrent networks as kernel methods
\cite{fermanian2021framing}, and through the signature kernel, which can be
computed by solving a Goursat PDE \cite{salvi2021goursat}. Sun et al.\
\cite{sun2024recursig} use signatures for explainable classification of
network time series. A survey of recent applications of rough path theory
is given by Fermanian et al.\ \cite{fermanian2023newdirections}.

\subsection{Koopman operator and dynamic mode decomposition}

Koopman \cite{koopman1931} showed that a nonlinear dynamical system can be
described by a linear operator acting on functions of its state. Mezi\'c
\cite{mezic2005spectral} used the spectrum of this operator for model
reduction and decomposition. Dynamic mode decomposition
\cite{schmid2010dmd} extracts dynamic modes and their eigenvalues directly
from data; the book by Kutz et al.\ \cite{kutz2016dmd} explains how these
approximate the Koopman spectrum and covers the variants of DMD, including
delay coordinates, which we use to apply DMD to a single series. Lusch et al.\ \cite{lusch2018deepkoopman} learn the
functions in which the dynamics become linear with a neural network. In
this work DMD serves as a very small, interpretable feature set; the case
study shows that eight eigenvalue features are not enough for
classification.

\subsection{Benchmarks and expected results}

Large comparisons of time series classifiers are made on public
archives with fixed train/test splits, such as the univariate study of
Bagnall et al.\ \cite{bagnall2017bakeoff}. For multivariate data, the UEA
archive \cite{bagnall2018uea} provides 30 datasets, including the four
used here, with benchmark results of nearest-neighbour classifiers.
Table~\ref{tab:published} compares these published results with our
feature methods on clean data. The archive reports 1-NN with the
Euclidean distance (ED) and with dynamic time warping applied to each
dimension separately (DTW$_I$) or to all dimensions together (DTW$_D$).

\begin{table}[ht]
\centering
\caption{Test accuracy on the original train/test splits without missing
points. Top: published 1-NN baselines (un-normalised data, Table~3 of
\cite{bagnall2018uea}). Bottom: our feature methods
(Table~\ref{tab:features}).}
\label{tab:published}
\begin{tabular}{lcccc}
\toprule
Method & BasicMotions & Libras & NATOPS & ArticularyWR \\
\midrule
1-NN ED \cite{bagnall2018uea} & 0.675 & 0.833 & 0.850 & 0.970 \\
1-NN DTW$_I$ \cite{bagnall2018uea} & 1.000 & 0.894 & 0.850 & 0.980 \\
1-NN DTW$_D$ \cite{bagnall2018uea} & 0.975 & 0.872 & 0.883 & 0.987 \\
\midrule
Summary stats (ours) & 1.000 & 0.683 & 0.833 & 0.893 \\
Signature $m = 3$ (ours) & 0.900 & 0.700 & 0.872 & 0.980 \\
\bottomrule
\end{tabular}
\end{table}

The 1-NN DTW classifiers are the bar that a new method should reach on
these datasets. The depth-3 signature with logistic regression comes
close on NATOPS and ArticularyWordRecognition (within 0.011 and 0.007 of
DTW$_D$) but is clearly below on Libras (0.700 against 0.894), where the
best published baseline is still far ahead. The target for clean data is
therefore at least the DTW level.
% Later: add stronger recent multivariate classifiers (e.g. the multivariate
% bake-off of Ruiz et al.) once their publication details and numbers are verified.
Under missing data, the result of Kidger et al.\ \cite{kidger2020ncde} on
CharacterTrajectories suggests that a Neural CDE can keep its accuracy
almost unchanged up to 70\% dropped points; this is the behaviour that
RQ2 tests on our datasets.
